\subsection{Mixed masses and concavity}
\label{sec:ordinary-mixed-mass-concavity}
We now allow the currents and their background classes to vary. A positive pure mass makes each class big by the preceding theorem. Solving relative equations against one smooth measure then converts the local mixed lower bound into a total-mass inequality, while the singularity envelopes keep the original mixed mass unchanged.

\begin{theorem}\label{thm:ordinary-mixed-mass-geometric-mean}
Let $X$ be connected, compact and K\"ahler, of dimension $n\geq1$. For closed positive $(1,1)$-currents $T_1,\ldots,T_n$,
\begin{equation}\label{eq:ordinary-mixed-mass-geometric-mean}
 \int_X\langle T_1\wedge\cdots\wedge T_n\rangle
       \geq\prod_{i=1}^n\left(\int_X\langle T_i^n\rangle\right)^{1/n}.
\end{equation}
Consequently the functions
\[
 T\longmapsto\left(\int_X\langle T^n\rangle\right)^{1/n},
 \qquad T\longmapsto\log\left(\int_X\langle T^n\rangle\right)
\]
are concave on the cone of closed positive $(1,1)$-currents. For the second function use $\log0=-\infty$.
\end{theorem}
\begin{proof}
All displayed masses are finite. If a pure mass vanishes, the right side of \eqref{eq:ordinary-mixed-mass-geometric-mean} is zero, and positivity proves the inequality. Otherwise let $m_i=\int_X\langle T_i^n\rangle>0$ and replace $T_i$ by $m_i^{-1/n}T_i$. Positive multilinearity reduces the proof to $m_i=1$ for all $i$.

By \cref{thm:positive-nonpluripolar-mass-bigness}, each normalized current has big class. The global $\ddc$ lemma writes it as $T_i=\theta_i+\ddc u_i$ with smooth closed $\theta_i$. Set $\phi_i=P_{\theta_i}[u_i]$. Mixed envelope mass preservation and the positive-mass model characterization, \cref{prop:p-envelope-calculus,thm:positive-mass-model-characterization}, make $\phi_i$ a model potential of mass one.

Choose the smooth probability measure $\nu=\omega^n/\int_X\omega^n$. It gives no mass to pluripolar sets, since local pole sets have zero volume. Apply \cref{thm:nonpluripolar-relative-prescribed-measure-existence} separately in each big class to obtain
\[
 w_i\in\mathcal E(X,\theta_i,\phi_i),\qquad
 S_i=\theta_i+\ddc w_i\geq0,\qquad
                 \langle S_i^n\rangle=\nu.
\]
The relative full-mass characterization gives
$P_{\theta_i}[w_i]=P_{\theta_i}[\phi_i]=\phi_i$.
Apply mixed envelope mass preservation first to the tuple $w_i$ and then to the tuple $u_i$. It gives
\[
 \int_X\langle S_1\wedge\cdots\wedge S_n\rangle
 =\int_X\left\langle\bigwedge_{i=1}^n
               (\theta_i+\ddc\phi_i)\right\rangle
 =\int_X\langle T_1\wedge\cdots\wedge T_n\rangle.
\]
Meanwhile \cref{prop:nonpluripolar-multifactor-common-measure}, with every density equal to one, gives
$\langle S_1\wedge\cdots\wedge S_n\rangle\geq\nu$.
Its total mass is at least one, proving the normalized inequality. Undoing the positive scalings proves \eqref{eq:ordinary-mixed-mass-geometric-mean}.

Write $M(T)=\int_X\langle T^n\rangle$, and take closed positive $T,S$ and $0<t<1$. Positive multilinearity gives
\[
 M((1-t)T+tS)=\sum_{k=0}^n\binom nk(1-t)^kt^{n-k}
                           \int_X\langle T^k\wedge S^{n-k}\rangle.
\]
For $1\leq k\leq n-1$, use the proved inequality with repeated factors. The endpoints are $M(T),M(S)$, including when either vanishes. The binomial formula yields
\[
 M((1-t)T+tS)\geq
              \big((1-t)M(T)^{1/n}+tM(S)^{1/n}\big)^n.
\]
Taking nonnegative roots proves the first concavity statement; $t=0,1$ are identities. If both endpoint masses are positive, the increasing concave logarithm gives
\[
 \log M((1-t)T+tS)
               \geq(1-t)\log M(T)+t\log M(S).
\]
If an endpoint mass is zero, the right side is $-\infty$ for $0<t<1$, and the inequality remains valid. Endpoint values of $t$ are interpreted directly without multiplying $0$ by $-\infty$. For $n=1$ the mixed inequality is an identity, the mass is additive, and the same logarithm argument applies.
\end{proof}
\index{non-pluripolar mass!geometric-mean inequality}
\index{non-pluripolar mass!concavity}

Connectedness is essential for the geometric-mean conclusion. On two disjoint compact K\"ahler components of dimension $n\geq2$, choose $T_1$ strictly positive on the first and zero on the second, and choose $T_2$ with the opposite support. Complete the tuple with currents strictly positive on both components. Every pure total mass is positive, but the mixed product vanishes. The theorem applies separately on each connected component.
